%%%
%
% $Autor: C. Crety, H. Mey, J. Perewersenko $
% $Datum: 2025-07-10 $

% !TeX spellcheck = de_DE/GB
% !TeX program = pdflatex
% !BIB program = biber/bibtex
% !TeX encoding = utf8
%
%%%



\chapter{Setup für den Arduino Nano 33 BLE Sense}

\section{Einführung}

In diesem Kapitel wird der Prozess beschrieben, wie der Arduino Nano 33 BLE Sense mit einem Computer verbunden und die notwendigen Einstellungen vorgenommen werden, um mit der Programmierung zu beginnen.

Behandelt werden:
\begin{itemize}
	\item die Installation erforderlicher Tools
	\item die die Konfiguration in der Arduino IDE
	\item ein einfacher Beispiel-Sketch zur Überprüfung der Funktion
\end{itemize}

Zur Überprüfung der Konfiguration kann einer der vordefinierten Beispiel-Sketches genutzt werden, z. B. der einfache \FILE{blnik.ino}, der mit der IDE ausgeliefert wird.

\section{Konfiguration für den Arduino Nano 33 BLE Sense}

Um den \textbf{Arduino Nano 33 BLE Sense} in einer offline Umgebung zu programmieren sind folgende Schritte notwendig:

\begin{enumerate}
	\subsection{Installation des Treibers}
	\item \textbf{Arduino IDE installieren:} Installieren Sie die aktuelle Version der Arduino IDE auf Ihrem Computer.
	\item \textbf{Installation der Boardpakete:} Öffnen Sie die \menu{Arduino IDE} und wählen Sie im Menü \menu{Tools}, in der oberen linken Ecke.
	\item Im Menü \menu{Tools}, wählen Sie den \menu{Board Manager}.
	\item Suchen Sie im Fenster \menu{Board Manager} nach \textbf{Arduino Nano 33 BLE Sense}.
	\item Wählen Sie in den Suchergebnissen das Paket \menu{Arduino Mbed OS Nano Boards} und clicken Sie \menu{Install}, um das Paket zu installieren.
	\item Schließen Sie das Board über einen USB-A zu Micro-USB Kabel an den Computer an.
\end{enumerate}

\begin{center}
	\begin{tikzpicture}
		\node (G) at (0,0) {\includegraphics[width=8cm]{Arduino/ArdiunoIDE2/boardmanager1.png}};
		
		\draw[red,domain=0:360,line width=1.4pt] plot ({-2.6+1*cos(\x)},{0.7
			+0.5*sin(\x)});
		
	
		
		
	\end{tikzpicture}
				\captionof{figure}{Arduino Mbed OS Nano Boards Installation}\label{fig:ArduinoMbedOSNanoBoardsInstallation}
\end{center}	



%%%%%%%%%%%%%%%%%%%%%%%%%%%%%%%%%%%%%%%%%%%%%%%%%%%%%%%%%%%%%%%%%%%%%%%%%%%%%%%%%%%%%%%

\subsection{Verbindung und Konfiguration des Arduino-Boards mit dem Computer}

%%%%%%%%%%%%%%%%%%%%%%%%%%%%%%%%%%%%%%%%%%%%%%%%%%%%%%%%%%%%%%%%%%%%%%%%%%%%%%%%%%%%%%%

Um entweder ein vordefiniertes Beispiel oder ein eigenes Programm auf dem Arduino Nano 33 BLE Sense auszuführen, befolgen Sie die folgenden Schritte::

\begin{enumerate}
	
	\item Verbinden Sie das Arduino-Board mit dem Computer.
	\item Öffnen Sie die Arduino IDE, dann erscheint ein leeres Projekt mir den Standardfunktionen \PYTHON{void setup()} und \PYTHON
	{void loop()}.
	\item Gehen Sie im Menü zu \menu{Tools}, wählen Sie \menu{Board} und wählen Sie das angeschlossene Board \textbf{Arduino Nano 33 BLE Sense}, wie in Abbildung~\ref{fig:SelecttheConnectedboard-hereArduinoNano33BLESense}.
	
\end{enumerate}

Damit ist die IDE so konfiguriert, dass sie das Board erkennt und mit ihm kommunizieren kann.

\begin{figure}[H]\centering
	\includegraphics[width=8cm]{Arduino/ArdiunoIDE2/SelectTheConnectedBoardArduinoNano33BLESense}
	\caption{Select the Connected board - here Arduino Nano 33 BLE Sense}
	\label{fig:SelecttheConnectedboard-hereArduinoNano33BLESense}		
\end{figure}

%%%%%%%%%%%%%%%%%%%%%%%%%%%%%%%%%%%%%%%%%%%%%%%%%%%%%%%%%%%%%%%%%%%%%%%%%%%%%%%%%%%%%%%

\subsection{Auswahl des passenden Ports}

%%%%%%%%%%%%%%%%%%%%%%%%%%%%%%%%%%%%%%%%%%%%%%%%%%%%%%%%%%%%%%%%%%%%%%%%%%%%%%%%%%%%%%%

Nach der Auswahl des Arduino Nano 33 BLE Sense, muss nun der richtige Port ausgewählt werden. Dafür versetzen Sie das Board in den Bootloader-Modus, drücken Sie dazu den weißen Reset-Knopf (siehe Abbildung ~\ref{fig:ArduinoNano33BLESenseResetButton}).

\begin{center}
	\begin{tikzpicture}
		\node (G) at (0,0) {\includegraphics[width=8cm]{Arduino/ArdiunoIDE2/resetbutton.png}};
		
		\draw[green,domain=0:360,line width=1.4pt] plot ({-1.5+0.4*cos(\x)},{-1.6
			+0.4*sin(\x)});
	\end{tikzpicture}
		\captionof{figure}{Arduino Nano 33 BLE Sense Reset Button}\label{fig:ArduinoNano33BLESenseResetButton}		
\end{center}


Wenn der Bootloader-Modus aktiviert ist, leuchtet die orangefarbene LED, wie in Abbildung ~\ref{fig:ArduinoNano33BLESenseResetButton}			



\begin{center}
	\begin{tikzpicture}
		\node (G) at (0,0) {\includegraphics[width=8cm]{Arduino/ArdiunoIDE2/resetbutton.png}};
		
		\draw[green,domain=0:360,line width=1.4pt] plot ({-2.8+0.4*cos(\x)},{-0.8
			+0.4*sin(\x)});
			
		\draw[orange,domain=0:360,line width=1.4pt] plot ({-2.8+0.4*cos(\x)},{-2.6
	+0.4*sin(\x)});
			
			
			
	\end{tikzpicture}
		\captionof{figure}{green and orange Builtin-LED}\label{fig:ArduinoNano33BLESenseResetButton}	
\end{center}	


Nun müssen Sie den Port auswählen. Gehen Sie dazu zu \menu{Tools}, wählen Sie \menu{Port} und versichern Sie sich, dass der richtig Port ausgewählt ist. ~\ref{fig:SelectAvailablePortforUploadingArduinoSketch}	



\begin{figure}[H]\centering
	\includegraphics[width=8cm]{Arduino/ArdiunoIDE2/SelectAvailablePortForUploadingArduinoSketch}
	\caption{Select Available Port for Uploading Arduino Sketch}
	\label{fig:SelectAvailablePortforUploadingArduinoSketch}		
\end{figure}




\subsection{Sketch hochladen und überprüfen}\label{uploadcode}

After ensuring that the appropriate port is selected, the next step is to upload the Arduino program. \textbf{Before} uploading the program, it is considered best practice to verify it first. This process will indicate whether any errors or warnings exist in the program. Once the program is successfully verified, it can be safely uploaded by clicking the button \menu{Upload} located below the file section, as shown in Figure ~\ref{fig:Upload the Program in Arduino board}




After selecting the board and port, confirm that the Arduino Nano 33 BLE Sense is properly connected to the Arduino IDE. This can be verified by checking the message in the bottom right corner of the IDE window, which should display the connected board and port.





\begin{figure}[H]\centering
	\includegraphics[width=8cm]{Arduino/ArdiunoIDE2/UploadTheProgramInArduinoBoard}
	\caption{Upload the Program in Arduino board}
	\label{fig:Upload the Program in Arduino board}		
\end{figure}




After uploading, the code will be compiled, and any issues in the program will be displayed in the bottom black window. Once the code is successfully uploaded and compiled on the Arduino board, it is necessary to reselect the port, as done previously. Navigate to the \menu{Tools menu}, select \menu{Port}, and ensure the correct port is selected, as shown in Figure ~\ref{fig:SettingthePort}, in order to view the output in the Serial Monitor.



\begin{figure}[H]\centering
	\includegraphics[width=8cm]{Arduino/ArdiunoIDE2/SettingThePort}
	\caption{Setting the Port}
	\label{fig:SettingthePort}		
\end{figure} 





\section{Test the Configuration}

In this subsection, we will test whether the board is properly connected to the Arduino IDE by running an Example Blink Sketch. This simple test will make an onboard LED light up, verifying both the connection and basic functionality of the Arduino Nano 33 BLE Sense. By the end of this section, you will have confirmed that the hardware and software are working together seamlessly.


\subsection{Testing the Steps by an Example Sketch \FILE{blink.ino}}


The Arduino IDE contains a set of built-in examples for testing purposes. To verify the configuration and set up the board, the basic example \FILE{blink.ino} can be opened first. After uploading the example \FILE{blink.ino} the Builtin LED blinks with 2 Hz. 


To begin, open the Arduino IDE on the computer and follow the steps below:

\begin{enumerate}



\item \textbf{Open the Examples:} Once it's open, click on the \menu{File} menu and hover over \menu{Examples} in the dropdown menu. 

\item \textbf{Selecting the Example Sketch:} A list of example categories will appear. Under the Built-in Examples, go to \menu{01.Basics} and click on \menu{Blink} as shown in ~\ref{fig:LED-ExampleTest}.

\item \textbf{Blink Example Sketch:} After following steps 1 and 2, the  Example Sketch \FILE{blink.ino} will open in a new window within the IDE as seen in Figure ~\ref{fig:LED-Built_Code}.

\item \textbf{Upload Example Sketch \FILE{blink.ino}:} Click the \menu{Upload} button (right arrow icon) in the Arduino IDE toolbar to compile and upload the example sketch to the Arduino Nano 33 BLE Sense board. Once the upload is complete, the orange built-in LED starts blinking. 



\end{enumerate}


\begin{center}
	\includegraphics[width=8cm]{Arduino/ArdiunoIDE2/blink_example_IDE2.3.3.png}
	\captionof{figure}{LED Example Sketch}\label{fig:LED-ExampleTest}		
\end{center} 



\begin{center}
	\includegraphics[width=8cm]{Arduino/ArdiunoIDE2/ledblink_code.png}
	\captionof{figure}{LED-Built Example}\label{fig:LED-Built_Code}		
\end{center} 



With the successful blinking of the built-in orange LED, it is confirmed that the Arduino Nano 33 BLE Sense is properly connected to the Arduino IDE and its basic functionality has been verified. The board is now ready for further development and projects.






%
\subsection{Description of Basic Sketch for Printing 'Hello'}

To test the development environment and the basic functionality of the hardware, a simple example sketch that controls only one LED is suitable. This sketch provides the Arduino \ac{ide}. Under the path \FILE{File/Examples/01.Basics}, small example sketches can be selected. Here, the example \FILE{Fade} is used. In this case, the built-in LED, which is an RGB LED, is utilized. 


\begin{lstlisting}
	int brightness = 0;  // how bright the LED is
	int fadeAmount = 5;  // how many points to fade the LED by
	
	// the setup routine runs once when you press reset:
	void setup() {
		// declare LED_BUILTIN to be an output:
		pinMode(LED_BUILTIN, OUTPUT);
	}
	
	// the loop routine runs over and over again forever:
	void loop() {
		// set the brightness of LED_BUILTIN:
		analogWrite(LED_BUILTIN, brightness);
		
		// change the brightness for next time through 
		//the loop:
		brightness = brightness + fadeAmount;
		
		// reverse the direction of the fading at the ends 
		//of the fade:
		if (brightness <= 0 || brightness >= 255) {
			fadeAmount = -fadeAmount;
		}
		// wait for 30 milliseconds to see the dimming 
		//effect
		delay(30);
	}
\end{lstlisting}


As a result, the brightness of the built-in LED should gradually fade in and out.

%

